\subsection{The bias boundary}

To find the endpoint, we compare the conditional density's slopes
at reflected points. Their ratio tells us how much weight the
reflected density needs to create a decrease.

\begin{proposition}[a variational equation for the bias boundary]
\label{prop:bias-boundary}
For $a\in(1/2,1)$ put $U_a=\{t>0:h_a'(t)>0\}$. With the infimum
over an empty set interpreted as $+\infty$,
\begin{equation}\label{eq:bias-boundary}
 b_a=\min\left\{\frac12,\quad
       \inf_{t\in U_a}\frac{h_a'(-t)}{h_a'(t)+h_a'(-t)}\right\}.
\end{equation}
In particular $b_a\to0$ as $a\uparrow1$. If the infimum is attained
at $t>0$ with $h_a'(t)>0$ and $b_a<1/2$, then, with $q=1-b_a$,
\[
 g_{a,q}'(-t)=g_{a,q}''(-t)=0.
\]
At this point the smaller peak disappears.
\end{proposition}

\begin{proof}
The proof of Theorem~\ref{thm:biased-limit-modes} shows that
$h_a'(-t)\ge0$ for $t\ge0$. For $q\ge1/2$ bimodality is equivalent
to a decrease somewhere on the negative half-line, hence to
$g_{a,q}'(-t)<0$ for some $t>0$. Here
\[
 g_{a,q}'(-t)=q h_a'(-t)-(1-q)h_a'(t).
\]
This expression can be negative only if $t\in U_a$, in which case
it is negative exactly when
$1-q>h_a'(-t)/(h_a'(t)+h_a'(-t))$. Taking the union over $t$ gives
\eqref{eq:bias-boundary}, with equality excluded as in the interval
theorem. At an interior minimizer the ratio has derivative zero.
Differentiating its numerator and denominator gives the two displayed
equations. The moment estimate in
Corollary~\ref{cor:bias-boundary-rate} gives $b_a\to0$.
\end{proof}

\begin{theorem}[regularity of the bias boundary]\label{thm:bias-continuity}
The function $a\mapsto b_a$ is continuous on $(1/2,1)$, is continuously
differentiable on $(a_*,1)$, and extends continuously to $b_1=0$.
It is strictly decreasing on some interval immediately above $a_*$.
For every $a>a_*$, we have $0<b_a<1/2$,
and the infimum in
\eqref{eq:bias-boundary} has a unique minimizer $t_a\in(0,\kappa_a)$.
This minimizer depends continuously on $a$ throughout $(a_*,1)$.
\end{theorem}

\begin{proof}
For $a>a_*$, the curvature partition in Theorem~\ref{thm:two-modes}
gives $g_a'(t)>0$ on $(0,\kappa_a)$ and $g_a'(t)<0$ on
$(\kappa_a,\infty)$. For $t\ge\kappa_a$ in $U_a$, the ratio in
\eqref{eq:bias-boundary} is therefore at least $1/2$.
Since $b_a<1/2$, it is enough to search $(0,\kappa_a)$
(we still have to show the infimum is attained).
Put
\[
 D_a(t)=h_a'(t)+h_a'(-t),\qquad
 R_a(t)=\frac{g_a'(t)}{D_a(t)}\quad(0<t<\kappa_a).
\]
Lemmas~\ref{lem:balanced-start} and~\ref{lem:positive-conditional-mode}
give $h_a'(-t)>0$ for every $t\ge0$.
We therefore have $D_a>0$, $0<R_a<1/2$, and
$1/2-b_a=\sup_{0<t<\kappa_a}R_a(t)$. Also $R_a(0+)=0$.

We first rule out a strict valley in the ratio. We claim that
\begin{equation}\label{eq:bias-quasiconcave}
 R_a(v)\ge\min\{R_a(u),R_a(w)\}
 \qquad(0<u<v<w<\kappa_a).
\end{equation}
Otherwise, choose $s\in(0,1/2)$ strictly between the middle value
and both outer values. With $q=1/2+s$,
$g_{a,q}'(-t)=D_a(t)(s-R_a(t))$ then has signs $-,+,-$
at $-w,-v,-u$, and is positive sufficiently close to zero on the left.
Positivity of the density
and its zero limits at infinity supply a positive derivative farther
left and a negative derivative farther right. There would be three
local maxima, contrary to Theorem~\ref{thm:biased-limit-modes}.

At $\kappa_a$, we have $D_a(\kappa_a)=2h_a'(-\kappa_a)>0$.
Thus $R_a$ extends continuously to $[0,\kappa_a]$, with value
zero at both endpoints. Its positive maximum is attained in the
interior. At that point $h_a'(-t)>0$, so the maximum is strictly
below $1/2$. Hence $b_a>0$. We need attainment here (pointwise strict
inequality on an open interval would not suffice).
Equation~\eqref{eq:bias-quasiconcave} and analyticity make this
maximum unique, since two maximizers would force a constant interval.
They also make $R_a$ strictly increasing before its maximizer and
strictly decreasing after it.

Lemma~\ref{lem:analytic-density} gives joint continuity of all the
derivatives used here. The signs of $g_a'$ on either side of its
unique positive zero give continuity of $\kappa_a$. Fix $a_0>a_*$
and choose
$0<l<t_{a_0}<r<\kappa_{a_0}$. The strict inequalities
$R_a(l),R_a(r)<R_a(t_{a_0})$ persist near $a_0$.
For $v<l$, apply \eqref{eq:bias-quasiconcave} to
$v,l,t_{a_0}$. Since $R_a(l)<R_a(t_{a_0})$, it gives
$R_a(v)\le R_a(l)$. Applying it to $t_{a_0},r,v$ gives
$R_a(v)\le R_a(r)$ for $v>r$.
Thus every nearby global maximum lies in the fixed interval $[l,r]$.
The ratios converge uniformly on this fixed compact interval,
so their maxima converge. Uniqueness gives convergence of the
maximizers too. This proves continuity above $a_*$.

At $a_*$, the curvature partition in Theorem~\ref{thm:two-modes}(a)
makes $g_{a_*}'$ strictly decrease from zero and then strictly increase
to zero. Thus $g_{a_*}'(r)<0$ for every $r>0$.
Joint continuity gives $g_a'(r)<0$ for all $a>a_*$ sufficiently close
to $a_*$, so $\kappa_a<r$. Since $r$ was arbitrary, $\kappa_a\to0$.
On the shrinking interval $[0,\kappa_a]$, joint continuity now gives
\[
 \sup_{0\le t\le\kappa_a}|g_a'(t)|\longrightarrow0,\qquad
 \inf_{0\le t\le\kappa_a}D_a(t)\longrightarrow2h_{a_*}'(0)>0,
\]
where the strict inequality is Lemma~\ref{lem:positive-conditional-mode}.
Consequently $1/2-b_a=\max R_a\to0$. Below the threshold $b_a=1/2$,
so this proves continuity at $a_*$ without the critical expansion.
Proposition~\ref{prop:bias-boundary} gives the limit at one.

For differentiability, the same three-point argument puts all nearby
maximizers in one compact interval. If $M(a)=1/2-b_a$, maximality gives,
for $s>0$,
\[
 R_{a+s}(t_a)-R_a(t_a)
 \le M(a+s)-M(a)
 \le R_{a+s}(t_{a+s})-R_a(t_{a+s}).
\]
Divide by $s$ and use the continuity of $t_a$ and
$\partial_aR_a$ from Lemma~\ref{lem:analytic-density}. The argument
for $s<0$ is the same. Thus
\begin{equation}\label{eq:bias-envelope-derivative}
 -b_a'=
 \frac{\partial_a h_a'(t_a)\,h_a'(-t_a)
       -h_a'(t_a)\,\partial_a h_a'(-t_a)}
      {\{h_a'(t_a)+h_a'(-t_a)\}^2}.
\end{equation}
The right side is continuous. No derivative of $t_a$ is needed.

Finally, put $C=\partial_a g_a''(0)|_{a=a_*}>0$ and
$H=h_{a_*}'(0)>0$, by Theorem~\ref{thm:threshold} and
Lemma~\ref{lem:positive-conditional-mode}.
Extend $R_a=g_a'/D_a$ to a fixed neighborhood of zero, where
$D_a>0$ for $a$ close to $a_*$. Evenness and
Lemma~\ref{lem:analytic-density} give, jointly as $(a,t)\to(a_*,0)$,
\[
 g_a'(t)=o(t),\qquad \partial_a g_a'(t)=(C+o(1))t,\qquad
 D_a(t)=2H+o(1),\qquad \partial_aD_a(t)=O(1).
\]
For the first two relations, integrate $g_a''$ and $\partial_a g_a''$
from zero to $t$, using $g_{a_*}''(0)=0$.
It follows that
\[
 \partial_aR_a(t)=\left(\frac{C}{2H}+o(1)\right)t>0
\]
for all sufficiently small $t>0$ and $a$ sufficiently close to $a_*$.
Since $0<t_a<\kappa_a\to0$, the envelope formula gives $b_a'<0$
on some interval immediately above the threshold.
\end{proof}

\begin{corollary}[a rate for the bias boundary near one]\label{cor:bias-boundary-rate}
Put
\[
 \rho_a=\frac{\Gamma(1+a)^2}{(2a-1)\Gamma(2a)}-1.
\]
If $\rho_a<1/12$, then
\[
 0\le b_a\le\frac{4\rho_a}{1-4\rho_a}.
\]
In particular, as $a\uparrow1$,
$b_a\le8(1-a)+O((1-a)^2)$.
\end{corollary}

\begin{proof}
The conditional limit has mean $\mu_a=1/\Gamma(1+a)>0$ and variance
$\rho_a\mu_a^2$, by the moment formulas of~\cite[Theorem~1.4]{glr2026}.
Consider the three adjacent intervals of length $\mu_a$ centered at
$-\mu_a$, $0$ and $\mu_a$. Chebyshev's inequality puts at least
$1-4\rho_a$ of the conditional law in the right interval and at most
$4\rho_a$ in the middle. For a start to the left, reflect these
intervals. After mixing, the left interval has mass at least
$(1-q)(1-4\rho_a)$ and the right at least $q(1-4\rho_a)$.
The middle still has mass at most $4\rho_a$.

For $4\rho_a/(1-4\rho_a)<q<1-4\rho_a/(1-4\rho_a)$, the middle
has less mass than either outer interval. All three intervals have
the same length. So, by continuity, we can choose a point in each
with the middle value strictly smaller than the other two.
This rules out unimodality. The two-mode bound in
Theorem~\ref{thm:biased-limit-modes} gives exactly two modes, and
therefore the stated bound on $b_a$. Finally, the logarithmic derivative at $a=1$
of $1+\rho_a$ is $2\psi(2)-2-2\psi(2)=-2$.
Thus $\rho_a=2(1-a)+O((1-a)^2)$, which gives the linear bound.
\end{proof}

\subsection{The density at the threshold}

How far apart are the two modes just after the transition?
We compute this from derivatives at the threshold. Write
\[
 C_* = \left.\partial_a g_a''(0)\right|_{a=a_*},\qquad
 B_*=-g_{a_*}^{(4)}(0),\qquad H_*=h_{a_*}'(0).
\]
These constants are positive by Theorem~\ref{thm:threshold} and
Lemmas~\ref{lem:analytic-quartic} and~\ref{lem:positive-conditional-mode}.
Table~\ref{tab:critical-constants} gives their numerical enclosures.
The coefficient $Q_*$ is defined in Proposition~\ref{prop:bias-cusp}.



\begin{proposition}[the critical shape]\label{prop:critical-shape}
The density at the threshold satisfies
\[
 g_{a_*}(u)=g_{a_*}(0)-\frac{B_*}{24}u^4+O(u^6)
 \qquad(u\to0).
\]
For $a>a_*$ let $\kappa_a>0$ be the positive mode of $g_a$. As $a\downarrow a_*$,
\begin{align*}
 \kappa_a&=\sqrt{\frac{6C_*}{B_*}}\sqrt{a-a_*}
                 +o\bigl(\sqrt{a-a_*}\bigr),\\
 g_a(\kappa_a)-g_a(0)&=\frac{3C_*^2}{2B_*}(a-a_*)^2
                 +o\bigl((a-a_*)^2\bigr).
\end{align*}
\end{proposition}

\begin{proof}
Lemma~\ref{lem:analytic-density} gives joint continuity of all spatial
derivatives and their first memory derivatives. Taylor remainders
through order four are therefore uniform on a neighborhood of $a_*$.
At criticality, evenness, $g_{a_*}''(0)=0$ and the analytic sign
$g_{a_*}^{(4)}(0)=-B_*<0$ from Lemma~\ref{lem:analytic-quartic}
give the first assertion.

Put $\delta=a-a_*>0$. Uniform Taylor expansion gives
\[
 \frac{g_a'(u)}u=C_*\delta+o(\delta)
                  -\frac{B_*+o(1)}6u^2+O(u^4).
\]
For $u=k\sqrt\delta$ this has the sign of $C_*-B_*k^2/6$
when $k$ is fixed away from $\sqrt{6C_*/B_*}$. The unique positive
mode lies between any two such choices with opposite signs, by
Theorem~\ref{thm:two-modes}. Squeezing gives the mode expansion.
Substitute it into
\[
 g_a(u)-g_a(0)=\tfrac12\bigl(C_*\delta+o(\delta)\bigr)u^2
                  -\tfrac1{24}\bigl(B_*+o(1)\bigr)u^4+O(u^6)
\]
to obtain the height expansion.
\end{proof}

\begin{proposition}[how the bias interval opens]\label{prop:bias-cusp}
With $C_*,B_*,H_*$ as defined above, put
\[
 Q_* = \frac{\sqrt{2}\,C_*^{3/2}}{3H_*\sqrt{B_*}},
 \qquad Q_*\approx9.867.
\]
As $a\downarrow a_*$,
\begin{equation}\label{eq:bias-cusp}
 \frac12-b_a=Q_*(a-a_*)^{3/2}+o\bigl((a-a_*)^{3/2}\bigr).
\end{equation}
The unique minimizer $t_a$ of
Theorem~\ref{thm:bias-continuity} satisfies
\[
 t_a=\sqrt{\frac{2C_*}{B_*}}\sqrt{a-a_*}
       +o\bigl(\sqrt{a-a_*}\bigr)
     =\frac{\kappa_a}{\sqrt3}+o\bigl(\sqrt{a-a_*}\bigr).
\]
At $q=1-b_a$, the stationary point is $-t_a$.
\end{proposition}

\begin{proof}
Lemma~\ref{lem:analytic-density} supplies the required joint regularity.
Theorem~\ref{thm:threshold} and
Lemma~\ref{lem:positive-conditional-mode} give $C_*>0$ and $H_*>0$.

Write $\delta=a-a_*$. By \eqref{eq:bias-boundary},
\[
 \frac12-b_a=\max\left\{0,\sup_{t\in U_a}
     \frac{g_a'(t)}{h_a'(t)+h_a'(-t)}\right\}.
\]
The fair-shape theorem makes the numerator nonpositive for $t\ge \kappa_a$.
Proposition~\ref{prop:critical-shape} gives
$\kappa_a=O(\sqrt\delta)$, and $H_*>0$ makes both conditional derivatives
positive on this shrinking interval. Uniform Taylor expansion gives,
for bounded $v\ge0$,
\[
 \frac{g_a'(v\sqrt\delta)}{h_a'(v\sqrt\delta)+h_a'(-v\sqrt\delta)}
 =\frac{\delta^{3/2}}{2H_*}
       \left(C_*v-\frac{B_*v^3}{6}\right)+o(\delta^{3/2}).
\]
The cubic is largest at the unique positive point
$v=\sqrt{2C_*/B_*}$. Maximizing on $[0,\kappa_a/\sqrt\delta]$, whose
right endpoint tends to $\sqrt{6C_*/B_*}$, proves \eqref{eq:bias-cusp}.
Since the ratio vanishes at both endpoints, its positive maximum
is attained inside $(0,\kappa_a)$. The convergence is uniform
and the limiting cubic has a unique maximum, so every scaled
maximizer tends to
$\sqrt{2C_*/B_*}$. This proves the assertion about $t_a$.


\end{proof}

The two powers come from the width and height of the same cubic.
Near zero the derivative is approximately $C_*(a-a_*)t-B_*t^3/6$.
Its positive part has width of order $\sqrt{a-a_*}$ and height of
order $(a-a_*)^{3/2}$. The width sets the distance between the modes.
The height sets how much bias removes the smaller peak.
